% Suggested final budget: 1.0--1.25 pages.
\section{Policy and Implementation}
\label{sec:method}
\subsection{Conditional action denoising}
For observation history $o$ and clean action sequence $a^{(0)}$, a standard
noise-prediction objective is
\begin{equation}
  \mathcal{L}(\theta)=
  \mathbb{E}_{o,a^{(0)},k,\epsilon}
  \left[\left\|\epsilon-
  \epsilon_\theta\!\left(
  \sqrt{\bar\alpha_k}a^{(0)}+
  \sqrt{1-\bar\alpha_k}\epsilon,k,o\right)
  \right\|_2^2\right],
  \label{eq:denoising}
\end{equation}
where $k$ denotes the diffusion timestep and $\epsilon\sim\mathcal{N}(0,I)$.
This describes the denoising formulation used by Diffusion
Policy~\citep{chi2023diffusion}.
\todo{Align the notation with the actual implementation. Describe observation,
prediction, and executed-action horizons; conditioning; padding; action
normalization and inverse normalization; and inference sampling.}

\subsection{Observation encoder and backbones}
\todo{Describe the implemented RGB/proprioception encoder, weight initialization,
image preprocessing, and feature fusion. Explain UNet, Transformer, and MLP
conditioning and architecture parameters. Attribute adapted external code and
state what was implemented or changed by the group.}

\begin{figure}[htbp]
  \centering
  \figplaceholder[32mm]{Expert generation $\rightarrow$ RGB/proprioception
  dataset $\rightarrow$ observation encoder and action denoiser $\rightarrow$
  receding-horizon control $\rightarrow$ held-out evaluation.}
  \caption{Placeholder for the implemented training and evaluation pipeline.
  Replace with a labeled diagram of actual observations, actions, and horizons.}
  \label{fig:pipeline}
\end{figure}

\subsection{Training and computing resources}
\todo{Summarize optimizer, learning rate and schedule, batch size, optimizer
steps, diffusion schedule, EMA, checkpoint policy, and hardware. Use the final
run configurations, not obsolete planning documents. State task-specific
exceptions and provide configuration/artifact links in the reproducibility
statement.}
